% !tp-example Basic usage
\documentclass{article}
\begin{document}
\wcTpPython[query=square(7)]{
  \wcScript[type=tp/python]|
    def square(x):
        return x * x
  |
}
\end{document}

% !tp-example Javascript API
\documentclass{article}
\begin{document}
\wcScript[type=module]|
import "/components/python/python.js";

const example = document.createDocumentFragment();

const element1 = document.createElement("tp-python");
element1.setAttribute("query", "square(7)");
element1.append(document.createTextNode("\n  "));

const element2 = document.createElement("script");
element2.setAttribute("type", "tp/python");
element2.textContent = "\n    def square(x):\n        return x * x\n  ";
element1.append(element2);
element1.append(document.createTextNode("\n"));
example.append(element1);

document.body.append(example);
|
\end{document}

% !tp-example Persistent session
\documentclass{article}
\begin{document}
\wcTpPython[query=next(counter)]{
  \wcScript[type=tp/python]|
    from itertools import count

    counter = count(1)
  |
}
\end{document}

% !tp-example Async functions and output
\documentclass{article}
\begin{document}
\wcTpPython[query=await answer()]{
  \wcScript[type=tp/python]|
    import asyncio

    async def answer():
        print("Computing the answer")
        await asyncio.sleep(0)
        return 42
  |
}
\end{document}

% !tp-example External Python file
\documentclass{article}
\begin{document}
\wcTpPython[src=/docs/components/python/examples/squares.py,query=square(9)]{}
\end{document}

% !tp-example DOM elements and click events
\documentclass{article}
\begin{document}
\wcScript[type=module]|
  import "/components/button/button.js";
  import "/components/python/python.js";
|

\wcTpPython[query=setup_dom()]{
  \wcScript[type=tp/python,filename=dom-demo.py]|
    from js import parent
    from pyodide.ffi import create_proxy

    handlers = []

    def setup_dom():
        host = parent.document.getElementById("python-dom-demo")
        host.replaceChildren()
        for handler in handlers:
            handler.destroy()
        handlers.clear()
        button = parent.document.createElement("tp-button")
        button.textContent = "Click me"
        message = parent.document.createElement("p")
        message.textContent = "Waiting for a click."
        clicks = 0

        def clicked(event):
            nonlocal clicks
            clicks += 1
            message.textContent = f"Python handled {clicks} click(s)."
            print("Clicks:", clicks)

        handler = create_proxy(clicked)
        handlers.append(handler)
        button.addEventListener("click", handler)
        host.append(button, message)
        return "Button ready"
  |
}

\wcDiv[id=python-dom-demo,aria-live=polite]{}
\end{document}

% !tp-example Errors and recovery
\documentclass{article}
\begin{document}
\rawhtml|<tp-python query="divide(12, 0)">
  <script type="tp/python" filename="errors.py">
    def divide(numerator, denominator):
        return numerator / denominator
  </script>
</tp-python>|
\end{document}

% !tp-example Presentation modes and automatic reload
\documentclass{article}
\begin{document}
Each interpreter has its own session. Default and lite show the normal controls; simple opens the editor and hides the header and welcome information. Its query buttons are hidden.


In default and lite, open Edit source. Run the prefilled query to obtain 10. Change 10 to 20 in the source, then click Run or press Return in the query field. In simple, use Return in the query field. Shift+Return works too.


The result becomes 20 and “> Program reloaded.” appears before the query. Run again without editing: the result remains 20, with no new reload message. Merely editing the source does not execute it.


\subsection*{Default}

\wcTpPython[mode=default,query=answer()]{
  \wcScript[type=tp/python,filename=answer.py]|
    def answer():
        return 10
  |
}

\subsection*{Lite}

\wcTpPython[mode=lite,query=answer()]{
  \wcScript[type=tp/python,filename=answer.py]|
    def answer():
        return 10
  |
}

\subsection*{Simple}

\wcTpPython[mode=simple,query=answer()]{
  \wcScript[type=tp/python,filename=answer.py]|
    def answer():
        return 10
  |
}
\end{document}

% !tp-example Execution trace drawer
\documentclass{article}
\begin{document}
\wcTpPython[tracer,filename=factorial.py,query=print(factorial(5))]{
  \wcScript[type=tp/python]|
    def factorial(n):
        if n <= 1:
            return 1
        return n * factorial(n - 1)
  |
}
\end{document}
